\section{Рабочий проект}

В процессе разработки были выделены следующие модули системы: \texttt{common}, \texttt{macro}, \texttt{compiler}, \texttt{optimizer}, \texttt{generator}, \texttt{assembler}.


\subsection{Модуль <<common>>}

Описание функций модуля \texttt{common} представлено в таблице \ref{common:table}.

\renewcommand{\arraystretch}{0.8} % уменьшение расстояний до сетки таблицы
\begin{xltabular}{\textwidth}{|p{2.5cm}|p{2.5cm}|X|}
	\caption{Описание функций, используемых в модуле \texttt{common}\label{common:table}}\\
	\hline \centrow \setlength{\baselineskip}{0.7\baselineskip} Название функции & \centrow Аргументы функции & \centrow Описание функции \\
	\hline \centrow 1 & \centrow 2 & \centrow 3 \\ \hline
	\endfirsthead
	\caption*{Продолжение таблицы \ref{common:table}}\\
	\hline \centrow 1 & \centrow 2 & \centrow 3 \\ \hline
	\finishhead
	comp-err & msg, \&rest other & Вызывает исключение \texttt{comp-err} с сообщением \texttt{msg} и опциональной дополнительной информацией об ошибке \texttt{other}. \\ \hline
	is-nary & args & Проверка на присутствие \texttt{\&rest} аргумента в \texttt{args}. \\ \hline
	mk/add-func & name list \&rest other & Макрос для генерации функции \texttt{name}, добавляющей информацию компиляции функций в список \texttt{list} с дополнительной информацией \texttt{other}. \\ \hline
	num-fix-args & list num & Определяет число фиксированных аргументов в списке \texttt{list} с начальным значением \texttt{num}. \\ \hline
	remove-rest & list & Удаляет символ \texttt{\&rest} из списка аргументов \texttt{list}. \\ \hline
	make-nary-args & count args & Помещает необязательные аргументы из списка \texttt{args}, в котором \texttt{count} обязательных аргументов, в ещё один список. \\ \hline
	search-symbol & list name & Поиск информации о функции или примитиве с названием \texttt{name} в списке \texttt{list}. \\ \hline
	correct-lambda & f & Проверка на то, что S-выражение \texttt{f} является корректным лямбда-выражением. \\ \hline
\end{xltabular}
\renewcommand{\arraystretch}{1.0} % восстановление сетки


\subsection{Модуль <<macro>>}

Описание функций модуля \texttt{macro} представлено в таблице \ref{macro:table}.

\renewcommand{\arraystretch}{0.8} % уменьшение расстояний до сетки таблицы
\begin{xltabular}{\textwidth}{|p{2.5cm}|p{2.5cm}|X|}
	\caption{Описание функций, используемых в модуле \texttt{macro}\label{macro:table}}\\
	\hline \centrow \setlength{\baselineskip}{0.7\baselineskip} Название функции & \centrow Аргументы функции & \centrow Описание функции \\
	\hline \centrow 1 & \centrow 2 & \centrow 3 \\ \hline
	\endfirsthead
	\caption*{Продолжение таблицы \ref{macro:table}}\\
	\hline \centrow 1 & \centrow 2 & \centrow 3 \\ \hline
	\finishhead
	add-fix-macro & name env arity args body & Сгенерированная макросом \texttt{mk/add-func} функция для сохранения информации о макросе с фиксированным числом аргументов \texttt{name} с окружением \texttt{env}, арностью \texttt{arity}, параметрами \texttt{args} и телом \texttt{body} в список \texttt{*fix-macros*}. \\ \hline
	add-nary-macro & name env arity args body & Сгенерированная макросом \texttt{mk/add-func} функция для сохранения информации о макросе с переменным числом аргументов \texttt{name} с окружением \texttt{env}, арностью \texttt{arity}, параметрами \texttt{args} и телом \texttt{body} в список \texttt{*nary-macros*}. \\ \hline
	make-macro & list & Сохранить информацию о макросе. \texttt{list} -- \texttt{(name args body)}, где \texttt{name} -- название макроса, \texttt{args} -- список параметров макроса, \texttt{body} -- тело макроса. \\ \hline
	subst & sym env & Подставить значение символа или глобальной переменной \texttt{sym} с окружением \texttt{env}. \\ \hline
	macro-eval & expr env & Рекурсивное раскрытие макроса \texttt{expr} с окружением \texttt{env}. \\ \hline
	macro-eval-backquote & expr env & Макровычисление квазицитирования \texttt{expr} с окружением \texttt{env}. \\ \hline
	macro-eval-if & expr env & Макровычисление условия \texttt{expr} с окружением \texttt{env}. \\ \hline
	macro-eval-args & args env & Макровычисление списка аргументов \texttt{args} с окружением \texttt{env}. \\ \hline
	macro-eval-progn & expr env & Макровычисление последовательности \texttt{expr} с окружением \texttt{env}. \\ \hline
	extend-macro-env & env args vals & Расширение окружения подстановки макросов \texttt{env} символами \texttt{args} и их значениями \texttt{vals}. \\ \hline
	macro-eval-let & args env & Макровычисление let-выражения \texttt{args} с временным расширением окружения \texttt{env}. \\ \hline
	macro-eval-app-func & args body vals env & Применение пользовательской функции с параметрами \texttt{args}, телом \texttt{body}, аргументами \texttt{vals} и окружением \texttt{env} внутри макроса. \\ \hline
	check-arguments & f type count args & Проверка на то, для функции или примитива \texttt{f} типа \texttt{type} список аргументов \texttt{args} имеет правильную длину в сравнении с числом фиксированных аргументов \texttt{count}. \\ \hline
	macroexpand & args vals body & Раскрытие макроса с параметрами \texttt{args}, аргументами \texttt{vals} и телом \texttt{body}. \\ \hline
	macro-eval-app & f args env & Макро применение функции, макроса или примитива \texttt{f} с аргументами \texttt{args} и окружением \texttt{env}. \\ \hline
	macro-eval-funcall & f args env & Макровычисление вызова функции \texttt{f} с аргументами \texttt{args} и окружением \texttt{env}. \\ \hline
	macro-eval-cond & list env & Макровычисление cond-выражения \texttt{list} с окружением \texttt{env}. \\ \hline
	macro-eval-setq & var expr env & Присвоение макро символу \texttt{var} значения \texttt{expr} с окружением \texttt{env}. \\ \hline
	macro-eval-function & s & Макровычисление функционального объекта \texttt{s}. \\ \hline
	macro-expand-progn & expr env & Макро раскрытие последовательности \texttt{expr} с окружением \texttt{env}. \\ \hline
\end{xltabular}
\renewcommand{\arraystretch}{1.0} % восстановление сетки


\subsection{Модуль <<compiler>>}

Описание функций модуля \texttt{compiler} представлено в таблице \ref{compiler:table}.

\renewcommand{\arraystretch}{0.8} % уменьшение расстояний до сетки таблицы
\begin{xltabular}{\textwidth}{|p{2.5cm}|p{2.5cm}|X|}
	\caption{Описание функций, используемых в модуле \texttt{compiler}\label{compiler:table}}\\
	\hline \centrow \setlength{\baselineskip}{0.7\baselineskip} Название функции & \centrow Аргументы функции & \centrow Описание функции \\
	\hline \centrow 1 & \centrow 2 & \centrow 3 \\ \hline
	\endfirsthead
	\caption*{Продолжение таблицы \ref{compiler:table}}\\
	\hline \centrow 1 & \centrow 2 & \centrow 3 \\ \hline
	\finishhead
	extend-env & env args & Расширяет окружение \texttt{env} новым кадром, состоящим из аргументов \texttt{args}. \\ \hline
	add-func & name env arity args body & Сгенерированная макросом \texttt{mk/add-func} функция для сохранения информации о пользовательских функциях с фиксированным числом аргументов \texttt{name} с окружением \texttt{env}, арностью \texttt{arity}, параметрами \texttt{args} и телом \texttt{body} в список \texttt{*fix-functions*}. \\ \hline
	add-nary-func & name env arity args body & Сгенерированная макросом \texttt{mk/add-func} функция для сохранения информации о пользовательских функциях с переменным числом аргументов \texttt{name} с окружением \texttt{env}, арностью \texttt{arity}, параметрами \texttt{args} и телом \texttt{body} в список \texttt{*nary-functions*}. \\ \hline
	add-local-func & name env arity args body & Сгенерированная макросом \texttt{mk/add-func} функция для сохранения информации о локальных функциях \texttt{name} с окружением \texttt{env}, арностью \texttt{arity} и меткой \texttt{real-name} в список \texttt{*local-functions*}. \\ \hline
	inner-compile & expr env tail & Рекурсивная функция компиляции выражения \texttt{expr} c окружением \texttt{env} в промежуточную форму. \texttt{tail} -- является ли текущее выражение хвостовым. \\ \hline
	compile-func-body & name args body env & Компиляция тела функции \texttt{name} с параметрами \texttt{args}, телом \texttt{body} и окружением \texttt{env}. \\ \hline
	compile-lambda & name args body env & Регистрирует функцию \texttt{name} и компилирует её тело \texttt{body}. \texttt{args} -- параметры функции, \texttt{env} -- окружение. \\ \hline
	find-func & f &  Определение типа функции \texttt{f} и поиск её информации в соответствующем списке. \\ \hline
	compile-progn & lst env tail & Компилирует блок \texttt{progn} со списком выражений \texttt{lst} и окружением \texttt{env}. \texttt{tail} -- является ли текущее выражение хвостовым. \\ \hline
	inc-function-ref & name & Увеличить число ссылок на функцию \texttt{name} \\ \hline
	compile-application & f args env tail & Компиляция применения функции \texttt{f} с аргументами \texttt{args} в окружении \texttt{env}. \texttt{tail} -- является ли текущее выражение хвостовым. \\ \hline
	compile-function & f env & Компиляция создания замыкания функции \texttt{f} с окружением \texttt{env}. \\ \hline
	compile-defun & expr env & Компилирует объявление функции \texttt{expr} с помощью \texttt{defun} в окружении \texttt{env}. \\ \hline
	add-global & sym & Добавляет глобальную переменную \texttt{sym} и возвращает её индекс в списке глобальных переменных \texttt{*global-variables*}. \\ \hline
	find-global-var & var & Производит поиск переменной в глобальном окружении по символу \texttt{var}. Возвращает индекс в списке глобального окружения, если символ найден, иначе nil. \\ \hline
	var-search & params var & Вычисление порядкового номера переменной \texttt{var} в списке \texttt{params} с учетом \texttt{\&rest}. \\ \hline
	find-local-var & var env & Производит поиск переменной в локальном окружении \texttt{env} по символу \texttt{var}. Возвращает индекс переменной в окружении, если символ найден, иначе \texttt{nil}. \\ \hline
	find-var & var env & Производит поиск переменной по символу \texttt{var} сначала в локальном окружении \texttt{env}, а затем в списке глобальных переменных \texttt{*global-variables*}. Возвращает список, состоящий из символа типа переменной (\texttt{GLOBAL}, \texttt{LOCAL} или \texttt{DEEP}) и индекса переменной. \\ \hline
	compile-setq & body env & Компилирует setq-выражение с телом \texttt{body} и окружением \texttt{env}. \\ \hline
	compile-if & expr env tail & Компилирует if-выражение с телом \texttt{expr} и окружением \texttt{env}. \texttt{tail} -- является ли текущее выражение хвостовым. \\ \hline
	compile-constant & c & Компиляция константы \texttt{c}. \\ \hline
	compile-labels & lst env tail & Компилирует блок \texttt{labels} с телом \texttt{lst} и окружением \texttt{env}. \texttt{tail} -- является ли текущее выражение хвостовым. \\ \hline
	compile-variable & v env & Компиляция переменной \texttt{v} в окружении \texttt{env}. \\ \hline
	compile-backquote & expr env & Компиляция квазицитирования с телом \texttt{expr} и окружением \texttt{env}. \\ \hline
	compile-tagbody & body env & Компиляция \texttt{tagbody} с телом \texttt{body} и окружением \texttt{env}. \\ \hline
	compile-catch & args env & Компиляция \texttt{catch} с телом \texttt{args} и окружением \texttt{env}. \\ \hline
	compile-throw & args env & Компиляция \texttt{throw} с телом \texttt{args} и окружением \texttt{env}. \\ \hline
	compile-defconst & body env tail & Компиляция \texttt{defconst} с телом \texttt{body} и окружением \texttt{env} и сохранение этой константы в хеше \texttt{*const-vals*}. \texttt{tail} -- является ли текущее выражение хвостовым. \\ \hline
	compile-apply & func args pack env tail & Компиляция применения функционального объекта \texttt{func} с аргументами \texttt{args} и окружением \texttt{env}. \texttt{pack} -- нужно ли собирать аргументы в список (для примитива \texttt{apply}), \texttt{tail} -- является ли текущее выражение хвостовым. \\ \hline
	compile-args & args env & Компиляция аргументов \texttt{args} в окружении \texttt{env}. \\ \hline
	compile & expr & Анализ S-выражения \texttt{expr} и преобразование в эквивалентное выражение, состоящее из более простых операций.
\end{xltabular}
\renewcommand{\arraystretch}{1.0} % восстановление сетки


\subsection{Модуль <<optimizer>>}

Описание функций модуля \texttt{optimizer} представлено в таблице \ref{optimizer:table}.

\renewcommand{\arraystretch}{0.8} % уменьшение расстояний до сетки таблицы
\begin{xltabular}{\textwidth}{|p{2.5cm}|p{2.5cm}|X|}
	\caption{Описание функций, используемых в модуле \texttt{optimizer}\label{optimizer:table}}\\
	\hline \centrow \setlength{\baselineskip}{0.7\baselineskip} Название функции & \centrow Аргументы функции & \centrow Описание функции \\
	\hline \centrow 1 & \centrow 2 & \centrow 3 \\ \hline
	\endfirsthead
	\caption*{Продолжение таблицы \ref{optimizer:table}}\\
	\hline \centrow 1 & \centrow 2 & \centrow 3 \\ \hline
	\finishhead
	accumulator-loading & tree & Удаляет избыточные формы загрузки значений в аккумулятор. \texttt{tree} -- форма \texttt{SEQ}. \\ \hline
	trivial-condition & alter & Если условие формы \texttt{ALTER} -- константа, то считает значение условия и сокращает форму \texttt{ALTER} до одной из её веток. \texttt{alter} -- форма \texttt{ALTER}. \\ \hline
	simplify-arithmetic & tree & Заменяет вызов арифметического примитива \texttt{tree} с неограниченным количеством аргументов на вложенные вызовы примитивов с фиксированным количеством аргументов. \texttt{tree} -- вызов примитива с переменным числом аргументов. \\ \hline
	constant-folding & tree & Свёртка и распространение констант. \texttt{tree} -- константа или вызов примитива. \\ \hline
	expand-seqs & tree & Разворачивает вложенные \texttt{SEQ}. \texttt{tree} -- форма \texttt{SEQ}. \\ \hline
	dead-code-elimination & tree & Убирает из формы \texttt{SEQ} неиспользуемый код (последовательность команд загрузки аккумулятора (кроме последней команды и команды перед \texttt{RETURN})). \\ \hline
	tail-call & tree & Оптимизация хвостовой рекурсии. \texttt{tree} -- форма \texttt{TAIL-CALL} или \texttt{TAIL-NCALL}. \\ \hline
	unused-functions & tree & Удаление неиспользуемых функций. \texttt{tree} -- форма \texttt{FUNC}. \\ \hline
	search-abs-tree & exp nodes & Возвращает истину если найдены хотя бы один узел из списка nodes. \\ \hline
	one-ref-args & body & Проверка на то, что в теле функции \texttt{body} к каждому аргументу только одно обращение. \\ \hline
	search-closure-deep-ref & exp & Поиск всех \texttt{FIX-CLOSURE} и проверка каждого тела замыкания на наличие \texttt{DEEP-REF}. \\ \hline
	beta-exp & exp args level & Выполнить подстановку тела функции \texttt{exp} с аргументами \texttt{args} и номером кадра активации \texttt{level}. \\ \hline
	make-nary-param & num args & Создать выражение для подстановки аргументов \texttt{args}, \texttt{num} -- число фиксированных аргументов. \\ \hline
	beta-expansion & call & Подстановка тела функции \texttt{call}, когда функция вызывается один раз и это возможно. \\ \hline
	can-inline & f & Возможно ли встраивание функции \texttt{f}. \\ \hline
	side-effects & exp & Есть ли в выражении \texttt{exp} побочные эффекты. \\ \hline
	optimize-tree1 & tree & Первый проход оптимизации промежуточной формы \texttt{tree}. Возвращает оптимизированное дерево. \\ \hline
	optimize-tree2 & tree & Второй проход оптимизатора -- встраивание функций, удаление после встраивания. \\ \hline
	optimize-tree & tree & Оптимизация промежуточной формы \texttt{tree}. \\ \hline
\end{xltabular}
\renewcommand{\arraystretch}{1.0} % восстановление сетки


\subsection{Модуль <<generator>>}

Описание функций модуля \texttt{generator} представлено в таблице \ref{generator:table}.

\renewcommand{\arraystretch}{0.8} % уменьшение расстояний до сетки таблицы
\begin{xltabular}{\textwidth}{|p{2.5cm}|p{2.5cm}|X|}
	\caption{Описание функций, используемых в модуле \texttt{generator}\label{generator:table}}\\
	\hline \centrow \setlength{\baselineskip}{0.7\baselineskip} Название функции & \centrow Аргументы функции & \centrow Описание функции \\
	\hline \centrow 1 & \centrow 2 & \centrow 3 \\ \hline
	\endfirsthead
	\caption*{Продолжение таблицы \ref{generator:table}}\\
	\hline \centrow 1 & \centrow 2 & \centrow 3 \\ \hline
	\finishhead
	emit & val & Добавить инструкцию \texttt{val} в массив линейных инструкций \texttt{*program*}. \\ \hline
	inner-generate & expr & Рекурсивная генерация кода для скомпилированного выражения \texttt{expr}. \\ \hline
	generate-if & expr & Генерация ветвления. \texttt{expr} -- список из условия, ветки по истине и ветки по лжи. \\ \hline
	generate-2-params & expr & Генерация функции \texttt{expr} с 2-мя параметрами. \\ \hline
	generate-deep & i j expr & Генерация формы \texttt{DEEP-SET}. \texttt{i} -- смещение в окружении, \texttt{j} -- индекс в кадре активации, \texttt{expr} -- присваиваемое выражение. \\ \hline
	generate-set & expr & Генерация форм \texttt{GLOBAL-SET} и \texttt{LOCAL-SET}. \\ \hline
	generate-args & args set rev & Генерация вычисления аргументов \texttt{args}. \texttt{set} -- тип инструкции для каждого аргумента (например, \texttt{PUSH}), \texttt{rev} -- в обратном порядке. \\ \hline
	generate-fix-prim & type args & Генерация вызова примитива с фиксированным числом аргументов \texttt{type} и аргументами \texttt{args}. \\ \hline
	generate-nary-prim & type num args & Генерация вызова примитива с переменным числом аргументов \texttt{type} и аргументами \texttt{args}. \texttt{num} -- число фиксированных аргументов. \\ \hline
	generate-reg-call & name fix-num env args & Генерация обычного вызова функции \texttt{name} с аргументами \texttt{args} и окружением \texttt{env}. \texttt{fix-num} -- число фиксированных аргументов для функций с переменным числом аргументов. \\ \hline
	generate-let & count args body & Генерация let-формы с телом \texttt{body} и аргументами \texttt{args}, расширение окружения \texttt{env}. \\ \hline
	generate-closure & op name env body & Генерация замыкания. \texttt{op} -- инструкция создания замыкания (\texttt{PRIM-CLOSURE, \texttt{NPRIM-CLOSURE} или \texttt{FIX-CLOSURE}}), \texttt{name} -- название функционального объекта, env -- окружение, \texttt{body} -- опциональное тело (для лямбда-выражения). \\ \hline
	generate-catch & tag body & Генерация формы \texttt{catch} с меткой \texttt{tag} и телом \texttt{body}. \\ \hline
	generate-throw & tag res & Генерация формы \texttt{throw} с меткой \texttt{tag} и значением \texttt{res}. \\ \hline
	generate-tail-call & name fix-num depth args & Генерация хвостового вызова функции \texttt{name} с числом фиксированных аргументов \texttt{fix-num}, глубиной вызова относительно входа в функцию \texttt{depth} и аргументами \texttt{args}. \\ \hline
	generate-nary & args & Генерация списка необязательных аргументов \texttt{args}. \\ \hline
	generate-apply & func args pack & Генерация применения функции \texttt{func} к аргументам \texttt{args}. \texttt{pack} -- нужно ли объединять аргументы в список (для примитива \texttt{apply}). \\ \hline
	generate & expr & Генерация кода для скомпилированного выражения \texttt{expr}. \\ \hline
\end{xltabular}
\renewcommand{\arraystretch}{1.0} % восстановление сетки


\subsection{Модуль <<assembler>>}

Описание функций модуля \texttt{assembler} представлено в таблице \ref{assembler:table}.

\renewcommand{\arraystretch}{0.8} % уменьшение расстояний до сетки таблицы
\begin{xltabular}{\textwidth}{|p{2.5cm}|p{2.5cm}|X|}
	\caption{Описание функций, используемых в модуле \texttt{assembler}\label{assembler:table}}\\
	\hline \centrow \setlength{\baselineskip}{0.7\baselineskip} Название функции & \centrow Аргументы функции & \centrow Описание функции \\
	\hline \centrow 1 & \centrow 2 & \centrow 3 \\ \hline
	\endfirsthead
	\caption*{Продолжение таблицы \ref{assembler:table}}\\
	\hline \centrow 1 & \centrow 2 & \centrow 3 \\ \hline
	\finishhead
	last-cdr & list & Найти последний \texttt{cdr} в списке \texttt{list}. \\ \hline
	asm-emit & \&rest bytes & Записать байты \texttt{bytes} в список \texttt{bytecode}. \\ \hline
	assemble-label & inst & Ассемблирование метки \texttt{inst}. \\ \hline
	assemble-const & inst & Ассемблирование инструкции загрузки константы \texttt{inst} в \texttt{ACC}. \\ \hline
	assemble-prim & inst & Ассемблирование инструкции вызова примитива \texttt{inst}. \\ \hline
	assemble-inst & inst jmp-labels & Ассемблирование инструкции общего вида \texttt{inst}. \\ \hline
	assemble-first-pass & program & Первый проход ассемблера (замена мнемоник инструкций соответствующими байтами и сбор информации о метках и переходах). \\ \hline
	assemble-second-pass & first-pass-res & Второй проход ассемблера (замена меток переходов на относительные адреса). \\ \hline
	assemble & program & Превращает список инструкций с метками \texttt{program} в байт-код. \\ \hline
\end{xltabular}
\renewcommand{\arraystretch}{1.0} % восстановление сетки


\begin{comment}
\subsection{Модульное тестирование оптимизатора}

Модульные тесты для фазы оптимизации представлены на рисунках \ref{unit:image}, \ref{unit2:image}.

\begin{figure}[H]
\lstinputlisting[language=Lisp]{optimizer-tests.lsp}
\caption{Модульные тесты фазы оптимизации}
\label{unit:image}
\end{figure}

\begin{figure}[H]
\lstinputlisting[language=Lisp]{optimizer-tests2.lsp}
\caption{Модульные тесты фазы оптимизации (продолжение)}
\label{unit2:image}
\end{figure}

\begin{figure}[H]
\begin{lstlisting}[language=Lisp]
(test-compile '((lambda (x) ((lambda (y) (cons x y)) 1)) 2)
'(FIX-LET 1 ((CONST 2)) (FIX-LET 1 ((CONST 1)) (PRIM CONS ((DEEP-REF 1 0) (LOCAL-REF 0))))))

(test-compile '(progn (defun fac (x) (if (equal x 1) 1 (* x (fac (- x 1))))) (fac 4))
'(SEQ (LABEL FAC (SEQ (ALTER (PRIM EQUAL ((LOCAL-REF 0) (CONST 1))) (CONST 1) (PRIM * ((LOCAL-REF 0) (REG-CALL FAC 0 ((PRIM - ((LOCAL-REF 0) (CONST 1)))))))) (RETURN))) (REG-CALL FAC 0 ((CONST 4)))))

(test-compile '(setq a #'(lambda (x y) (+ x y)))
'(GLOBAL-SET 2 (FIX-CLOSURE G598 (LABEL G598 (SEQ (PRIM + ((LOCAL-REF 0) (LOCAL-REF 1))) (RETURN))))))

(test-compile '(progn (defun test () 1) (setq a #'test))
'(SEQ (LABEL TEST (SEQ (CONST 1) (RETURN))) (GLOBAL-SET 2 (FIX-CLOSURE TEST ()))))

(test-compile '(defun test () #'(lambda (x) x))
'(LABEL TEST (SEQ (FIX-CLOSURE G711 (LABEL G711 (SEQ (LOCAL-REF 0) (RETURN)))) (RETURN))))

(test-compile '(progn (defmacro test (x y) `(+ ,x ,y)) (test 1 2))
'(SEQ (NOP) (PRIM + ((CONST 1) (CONST 2)))))

(test-compile '(progn (setq a 1 b 2) `(a (,a) b (,b)))
'(SEQ (GLOBAL-SET 2 (CONST 1))
(SEQ (GLOBAL-SET 3 (CONST 2))
(PRIM CONS ((CONST A)
(PRIM CONS ((PRIM CONS ((GLOBAL-REF 2)
(CONST NIL)))
(PRIM CONS ((CONST B)
(PRIM CONS ((PRIM CONS ((GLOBAL-REF 3)
(CONST NIL)))
(CONST NIL))))))))))))
\end{lstlisting}  
\caption{Модульные тесты фазы компиляции (продолжение)}
\end{figure}

\subsection{Системное тестирование разработанного компилятора}

На рисунках \ref{main:image}, \ref{main2:image} и \ref{main3:image} представлен системный тест компилятора, покрывающий всё реализованное подмножество Common Lisp, на рисунках \ref{main4:image}, \ref{main5:image} и \ref{main6:image} -- результат этого теста.
\begin{figure}[H]
\begin{lstlisting}[language=Lisp]
;; *test-failed* - флаг - провалился ли хотя бы один тест.
(defvar *test-failed*)
;; *started-tests-count* - количество запущенных тестов.
(defvar *test-compile-failed*) ;; TMP
(defvar *started-tests-count* 0)
;; *started-tests-count* - количество успешных тестов.
(defvar *ok-tests-count* 0)


;; Тест компиляции, ассемблирования и выполнения программы.
(defun test (expr expected-res)
  (unless (= *started-tests-count* *ok-tests-count*)
    (setq *test-failed* t))
  (incf *started-tests-count*)
  (unless (or *test-failed*
              *test-compile-failed*)
    (let* ((tree (compile expr))
           (asmcode (generate tree))
           (bytecode (assemble asmcode)))
      (print "Compiler")
      (labels ((all-atom (lst)
                 (if (atom lst)
                     t
                     (if (atom (car lst))
                         (all-atom (cdr lst))
                         nil)))
               (format-tree (tree indent lparen)
                 (if (atom tree)
                     (concat (make-string indent #\ )
                             (if (> lparen 0) (make-string lparen #\() "")
                             (cond
                               ((symbolp tree) (symbol-name tree))
                               ((integerp tree) (inttostr tree))
                               (t (error "Unreachable"))))
                     (let ((res nil))
                       (setq res (format-tree (car tree) indent (++ lparen)))
                       (if (all-atom (cdr tree))
                           (app #'(lambda (elem)
                                    (setq res (concat res " "
                                      (cond
                                        ((symbolp elem) (symbol-name elem))
                                        ((integerp elem) (inttostr elem))
                                        (t (error "Unreachable"))))))
                                (cdr tree))
                           (app #'(lambda (elem)
                                    (setq res (concat res "\n" (format-tree elem (++ indent) 0))))
                                (cdr tree)))
                       (setq res (concat res ")"))
                       res))))
        (print (concat "\n" (format-tree tree 0 0) "\n")))
\end{lstlisting}
\caption{Функции для системных тестов}
\label{main:image}
\end{figure}

\begin{figure}[H]
\begin{lstlisting}[language=Lisp]
      (print "Generator")
      (dolist (inst asmcode)
        (print inst))
      (print "Assembler")
      (print bytecode)
      (print "Execution")
      (let ((res (vm-run bytecode)))
        (print res)
        (setq res (assert res expected-res))
        (print res)
        (if (eq (car res) 'fail)
            (setq *test-failed* t)
            (progn
              (incf *ok-tests-count*)
              (print '---------------------)))))))

;; Проверка, все ли тесты успешны
(defun check-tests ()
  (when (and (= *started-tests-count* *ok-tests-count*)
             (null *test-failed*)
             (null *test-compile-failed*))
    (print "All tests OK")))
\end{lstlisting}
\caption{Функции для системных тестов (продолжение)}
\label{main2:image}
\end{figure}

\begin{figure}[H]
\begin{lstlisting}[language=Lisp]
(test
 '(progn
   (setq y 3)
   (defmacro -- (x) `(- ,x 1))
   (defun f (x)
     (if (> x 1)
         (* x (f (-- x)))
         1))
   ((lambda (x) (f x)) y))
 6)

(check-tests)
\end{lstlisting}
\caption{Системный тест компилятора}
\label{main3:image}
\end{figure}

\begin{figure}[H]
\begin{lstlisting}[language=Lisp]
"Compiler"
"
(SEQ
 (GLOBAL-SET
  2
  (CONST 3))
 (SEQ
  (NOP)
  (SEQ
   (LABEL
    F
    (SEQ
     (ALTER
      (PRIM
       >
       ((LOCAL-REF 0)
        (CONST 1)))
      (PRIM
       *
       ((LOCAL-REF 0)
        (REG-CALL
         F
         0
         ((PRIM
          -
          ((LOCAL-REF 0)
           (CONST 1)))))))
      (CONST 1))
     (RETURN)))
   (FIX-LET
    1
    ((GLOBAL-REF 2))
    (REG-CALL
     F
     0
     ((LOCAL-REF 0)))))))
"
\end{lstlisting}
\caption{Результат фазы компиляции системного теста компилятора}
\label{main4:image}
\end{figure}

\begin{figure}[H]
\begin{lstlisting}[language=Lisp]
"Generator"
(CONST 3)
(GLOBAL-SET 2)
(JMP G41)
(LABEL F)
(LOCAL-REF 0)
(PUSH)
(CONST 1)
(PUSH)
(PRIM >)
(JNT G54)
(LOCAL-REF 0)
(PUSH)
(LOCAL-REF 0)
(PUSH)
(CONST 1)
(PUSH)
(PRIM -)
(PUSH)
(SAVE-ENV)
(SET-ENV 0)
(ALLOC 1)
(REG-CALL F)
(RESTORE-ENV)
(PUSH)
(PRIM *)
(JMP G55)
(LABEL G54)
(CONST 1)
(LABEL G55)
(RETURN)
(LABEL G41)
(GLOBAL-REF 2)
(PUSH)
(SAVE-ENV)
(ALLOC 1)
(LOCAL-REF 0)
(PUSH)
(SAVE-ENV)
(SET-ENV 0)
(ALLOC 1)
(REG-CALL F)
(RESTORE-ENV)
(RESTORE-ENV)
\end{lstlisting}
\caption{Результат фазы генерации системного теста компилятора}
\label{main5:image}
\end{figure}

\begin{figure}[H]
\begin{lstlisting}[language=Lisp]
"Assembler"
#(0 3 5 2 1 40 6 0 10 0 1 10 18 13 2 27 6 0 10 6 0 10 0 1 10 18 4 10 14 15 0 3 1 11 -27 16 10 18 5 1 4 0 1 12 4 2 10 14 3 1 6 0 10 14 15 0 3 1 11 -52 16 16)
"Execution"
6
(OK 6 6)
---------------------
"All tests OK"
\end{lstlisting}
\caption{Результат фазы ассемблирования системного теста компилятора и выполнения полученного байт-кода на целевой машине}
\label{main6:image}
\end{figure}

\end{comment}